\chapter{From CSP Roles to Multiparty Grammars}
\label{ch:csp-exercises}

\index{CSP}
\index{protocol role}
This part revisits examples from \citet{ryan2001modelling} and its official Casper
companion scripts. The goal is not to transliterate CSP operators. It is to use
well-known security protocols as demanding exercises for the grammar-based method:
multiple roles, fresh values, encrypted components, an active intruder, calculated
knowledge sets, agreement, and messages whose legality depends on the preceding trace.

\section{Can actors really be specified separately?}

Yes, provided composition is defined as synchronization over the same observable
events rather than as concatenation of independent scripts. Each actor accepts a
language over its own projection of the global event alphabet.

\begin{lstlisting}
structure RoleSpec where
  actor    : Actor
  accepts  : Grammar (ProjectedEvent actor)
  initialKnowledge : Knowledge actor

def composeRoles (roles : List RoleSpec) : GlobalGrammar :=
  synchronize roles where
    sendReceiveAgree := sameEventId && sameBytes && oppositeEndpoints
\end{lstlisting}

For a global trace $t$, let $t\!\upharpoonright_a$ retain events sent or received by
actor $a$. The composition accepts exactly when every local role accepts its view:

\[
  L(\operatorname{composeRoles}(R_1,\ldots,R_n))
  = \{t \mid \forall i,\; t\!\upharpoonright_{a_i} \in L(R_i)\}.
\]

One global message is therefore not duplicated into unrelated send and receive
events. It has a single identity, byte string, sender, and recipient; both endpoint
grammars constrain it. An unmatched send, an unexpected receive, or disagreement
about bytes makes the intersection empty at that trace prefix.

\subsection{Local state without global implementation state}

A role may bind parameters from the events it observes. Those bindings are a
parser environment for the trace, not a claim about implementation memory. A role
can say ``the nonce in this reply equals the nonce I previously sent'' without
specifying a table, variable layout, or retry mechanism.

\section{A transport-neutral byte algebra}

\index{bytes}
Security claims need more than message names. They need an algebra that says which
bytes expose which values. We still need not choose HTTP, protobuf, ASN.1, or a
particular cipher suite.

\begin{lstlisting}
abbrev Bytes := ByteArray

inductive WireTerm where
  | atom    : Name -> Bytes -> WireTerm
  | tuple   : List WireTerm -> WireTerm
  | encrypt : KeyRef -> WireTerm -> WireTerm
  | sign    : KeyRef -> WireTerm -> WireTerm

structure WireEvent where
  clock : ClockTimestamp
  id     : EventId
  priorIds : List EventId
  session : SessionId
  task   : TaskId
  from   : Actor
  to     : Actor
  term   : WireTerm
  bytes  : Bytes
  encodedBy : bytes = encodeCanonical term
\end{lstlisting}

At this level, encryption is a symbolic constructor with an encoding into bytes.
The intruder may concatenate, split visible tuples, encrypt with known keys, and
decrypt only with a derivable inverse key. A concrete refinement must replace this
idealized algebra with a transport and cryptographic implementation whose security
assumptions justify the abstraction.

\section{Yahalom as three separate actors}

\index{Yahalom protocol}
Yahalom is an excellent first exercise because an initiator, responder, and trusted
server exchange fresh nonces and receive a session key. The companion model states
secrecy for the session key and agreement claims involving the nonces and key.

Let $K_A$ and $K_B$ be long-term keys shared by the server with Alice and Bob,
$N_A$ and $N_B$ fresh nonces, and $K_{AB}$ the proposed session key. The following
actor grammars use symbolic wire terms whose canonical encodings are byte strings.

\subsection{Alice, the initiator}

\begin{lstlisting}
def aliceRole (alice bob server : Actor) (ka : SharedKey) : RoleSpec :=
  role alice "Yahalom initiator" do
    let na <- fresh Nonce "na"

    from_ alice; to_ bob
    send (atom "initiator nonce" na.bytes)

    from_ server; to_ alice
    receive (encrypt ka.ref
      (tuple [actor bob, sessionKey "kab", nonce na, nonce "nb"]))
      where_ field "na" == na
    let kab <- take SessionKey "kab"
    let nb  <- take Nonce "nb"
    let ticketForBob <- take WireTerm "ticketForBob"

    from_ alice; to_ bob
    send ticketForBob
    send (encrypt kab.ref (nonce nb))
\end{lstlisting}

Alice's \texttt{where\_} clause is the context-sensitive step: the received nonce
must equal the value generated earlier. The surface structure is context-free, but
legality depends on the binding environment accumulated from the prefix.

\subsection{Bob, the responder}

\begin{lstlisting}
def bobRole (alice bob server : Actor) (kb : SharedKey) : RoleSpec :=
  role bob "Yahalom responder" do
    from_ alice; to_ bob
    receive (atom "initiator nonce" ?naBytes)
    let na <- take Nonce "na"
    let nb <- fresh Nonce "nb"

    from_ bob; to_ server
    send (encrypt kb.ref
      (tuple [actor alice, nonce na, nonce nb]))

    from_ alice; to_ bob
    receive (encrypt kb.ref (tuple [actor alice, sessionKey "kab"]))
    let kab <- take SessionKey "kab"
    receive (encrypt kab.ref (nonce nb)) where_ field "nb" == nb
\end{lstlisting}

Bob need not know Alice's local control state. He constrains only the events visible
at his boundary and the values learned from them.

\subsection{The trusted server}

\begin{lstlisting}
def serverRole (alice bob server : Actor)
    (ka kb : SharedKey) : RoleSpec :=
  role server "Yahalom key server" do
    from_ bob; to_ server
    receive (encrypt kb.ref
      (tuple [actor alice, nonce "na", nonce "nb"]))
    let na  <- take Nonce "na"
    let nb  <- take Nonce "nb"
    let kab <- fresh SessionKey "kab"

    let bobTicket := encrypt kb.ref
      (tuple [actor alice, sessionKey kab])
    from_ server; to_ alice
    sendMany
      [ encrypt ka.ref
          (tuple [actor bob, sessionKey kab, nonce na, nonce nb])
      , bobTicket
      ]
\end{lstlisting}

\subsection{Global assembly and context conditions}

\begin{lstlisting}
def yahalom : GlobalGrammar :=
  composeRoles
    [ aliceRole .alice .bob .server keyAlice
    , bobRole .alice .bob .server keyBob
    , serverRole .alice .bob .server keyAlice keyBob
    , intruderRole .ivo
    ]
  |> where_ (distinct [.alice, .bob, .server, .ivo])
  |> where_ (freshValuesNeverRepeat ["na", "nb", "kab"])
  |> where_ (longTermKeysKnownOnlyToOwnersAndServer)
\end{lstlisting}

The three honest roles can be authored and checked separately. Composition exposes
incompatibilities: for example, if Alice expects the server's first encrypted tuple
to name Bob but the server grammar could name another responder, the shared global
event cannot satisfy both projections.

\section{The attacker is another grammar plus deduction}

\index{intruder model}
An active network attacker observes public wire events, blocks them, replays known
bytes, and emits any term derivable from current knowledge. It does not learn the
plaintext of an encryption merely because its encoded bytes are observable.

\begin{lstlisting}
inductive Derives : Knowledge -> WireTerm -> Prop where
  | initial : term in knowledge -> Derives knowledge term
  | tuple   : (forall x in xs, Derives knowledge x) ->
              Derives knowledge (.tuple xs)
  | project : Derives knowledge (.tuple xs) -> x in xs ->
              Derives knowledge x
  | encrypt : Derives knowledge key -> Derives knowledge body ->
              Derives knowledge (.encrypt key body)
  | decrypt : Derives knowledge (.encrypt key body) ->
              Derives knowledge (inverse key) -> Derives knowledge body

def intruderRole (ivo : Actor) : RoleSpec :=
  adversary ivo do
    observe publicWireEvents
    mayBlock
    mayReplay
    maySend fun knowledge term => Derives knowledge term
\end{lstlisting}

This is the grammar analogue of the CSP/Dolev--Yao attacker: possible hostile
messages are generated from observable history and constrained by a
context-sensitive derivability relation.

\section{Knowledge and agreement become theorems}

\index{secrecy}
\index{agreement}
Knowledge is calculated for every actor from every accepted trace, not merely the
happy-path scenario. A traditional secrecy claim is recovered by comparing that
set with the actors permitted to know the value. For Yahalom:

\begin{lstlisting}
theorem yahalom_session_key_knowers
    (trace : Trace WireEvent)
    (accepted : yahalom.Accepts trace)
    (honest : Uncompromised [.alice, .bob, .server]) :
    knowers (sessionKey accepted.kab) trace = [.alice, .bob, .server] := by
  -- proof uses the fixed symbolic-crypto library
  exact knowledgeClosureFromKeyOwnership accepted honest

theorem yahalom_bob_agrees_with_alice
    (trace : Trace WireEvent)
    (accepted : yahalom.Accepts trace)
    (bobDone : completed trace .bob) :
    exists run, peer run == .alice &&
      run.na == bobDone.na && run.nb == bobDone.nb := by
  exact agreementFromNonceChallenge accepted bobDone
\end{lstlisting}

The exact proof scripts depend on the library still to be built. The important
design constraint is already visible: theorem statements refer only to accepted
event traces, knowledge derived from those traces, and explicit compromise
assumptions.

\section{Otway--Rees: guards are not optional decoration}

\index{Otway--Rees protocol}
The Otway--Rees companion model introduces a run marker $M$ and guards that the two
participants differ. Each participant encrypts its nonce, marker, and actor names
under its long-term server key. The server returns separate encrypted components.

This example adds three useful \texttt{where\_} obligations:

\begin{lstlisting}
where_ alice != bob
where_ response.field "marker" == request.field "marker"
where_ encryptedPart.names == [alice, bob]
\end{lstlisting}

The grammar shape remains a small sequence of messages. The accepted language is
not context-free once equality of an unbounded marker and actor parameters matters.
Keeping these tests as named guards preserves the readable protocol skeleton while
making the extra context explicit.

The example also asks whether an observer can distinguish executions after one
candidate session key is substituted for another. That cannot be answered by a
single knowledge set: it is a relational property over pairs of trace
distributions, beyond a single CTL state formula, and should be represented as a
separate theorem family.

\section{Needham--Schroeder: public-key lookup as behaviour}

\index{Needham--Schroeder protocol}
The seven-message public-key example makes key discovery observable. Alice requests
Bob's public key from a server, receives a signed binding, sends an encrypted nonce
and identity to Bob, and later answers Bob's nonce. Bob performs the symmetric key
lookup for Alice.

Actor separation is particularly valuable here:

\begin{itemize}
  \item the key server grammar promises signed actor--public-key bindings;
  \item Alice's grammar accepts Bob's key only where the signature verifies and the
        named actor is Bob;
  \item Bob imposes the corresponding constraint for Alice;
  \item the nonce-response messages reuse previously bound values;
  \item the intruder grammar may substitute any derivable public key or replay any
        signed binding, but cannot manufacture a server signature.
\end{itemize}

\begin{lstlisting}
receive (sign serverSigningKey (tuple [actor bob, publicKey "pkb"]))
  where_ verifiesWith serverPublicKey
  where_ field "actor" == bob

send (encrypt pkb (tuple [nonce na, actor alice]))

receive (encrypt aliceSecretView (tuple [nonce na, nonce "nb"]))
  where_ field "na" == na
\end{lstlisting}

This demonstrates why ``all behaviour is observable'' must not mean ``all
plaintext is public.'' The globally recorded event exposes endpoints, timing,
kind, and wire bytes. Which symbolic values an actor can recover from those bytes
is governed by the deduction relation.

\section{What these exercises demand from the library}

The three case studies produce a concrete implementation checklist:

\begin{enumerate}
  \item actor-local grammar definitions and trace projection;
  \item synchronization by common event identity and bytes;
  \item typed symbolic variables, exports, imports, and prefix-dependent guards;
  \item a transport-neutral term-to-byte relation;
  \item explicit initial knowledge and compromise assumptions;
  \item attacker closure under symbolic construction and decryption;
  \item per-value knowledge sets, agreement, and relational indistinguishability
        theorem schemas;
  \item counterexamples rendered simultaneously as a global interaction diagram
        and as each actor's local projection.
\end{enumerate}

These are not security-specific embellishments. They are a rigorous test of whether
multiparty context-sensitive grammars can remain modular while expressing the
information flow on which real behavioural claims depend.
